\section{Introduction}
\IEEEPARstart{A}{traffic} signal controller that works at one junction may make poor decisions at another. The number and ordering of phases can change, downstream roads may have different capacities, and the same queue length can call for a different action under a different demand pattern. These changes make cross-network deployment a more demanding test than evaluating a controller on new random seeds of the network used for training. For an operator, the practical question is whether experience from other networks improves decisions at the new location after a limited amount of local preparation.

Many signal-control methods learn a policy or value function from interaction with a simulator. Others predict future traffic and choose an action from those predictions. In either case, the decision depends on how the available actions compare. Predicting the absolute traffic state accurately is useful only to the extent that the prediction distinguishes the consequences of competing phases. An error shared by every candidate can be large without changing the preferred action; a small error in the difference between two candidates can reverse the decision.

This observation motivates a narrower prediction task. We restore the same simulated traffic state, try each feasible phase, and measure the difference in the resulting short-horizon costs. The learner then predicts these within-state differences. This construction removes cost components common to all candidates and avoids assigning a fixed output coordinate to each phase identity. It does not remove errors that affect actions differently, and it does not establish that a short-horizon ranking will remain useful after the controller changes the traffic state through repeated decisions.

We investigate this approach through Causal-Factored Cross-City Model Transfer (\method). Each source-group component combines two supervised predictors: a restricted predictor that compares each candidate with a pressure-based reference, and a pairwise predictor that compares every feasible pair. The second predictor is constrained so that reversing the order of a comparison reverses its sign. Target simulation data select how much weight to give each predictor. The fitted controller subsequently scores feasible phases from the observed traffic state, without online model updates or counterfactual simulation at each decision.

The study addresses three questions. First, can this representation improve phase selection across networks relative to a restricted predictor and an unrestricted residual baseline? Second, do improvements in offline action ranking translate into lower waiting time in closed-loop operation? Third, does the use of source-network examples improve on training from target data, once the target observations are held fixed? These questions require different comparisons. A model can outperform one residual estimator while remaining worse than a simple pressure rule, and a larger source-aware model can outperform a smaller target model without proving a benefit from source information itself.

We develop the method on 18 networks in seven groups and evaluate it on external Los Angeles and Jinan benchmark networks. The final controller selects source-group components or a target anchor using target development runs, then freezes both its predictor and execution guard. The main comparison tests this frozen source selector on 56 paired seeds. Earlier models and their separate validation experiments are reported in Appendix~\ref{app:development-results}.

The contributions are threefold. We formulate transfer around matched phase-cost differences and give the associated invariance and ranking properties. We implement a restricted predictor, an antisymmetric pairwise predictor, and a source-selection procedure that chooses a fitted component or target anchor before deployment. We then evaluate both their control performance and the extent to which gains can be attributed to source data, using pressure baselines and additional seeds on both targets. The results distinguish improvements in phase-cost prediction from improvements in sustained traffic control.
